\documentclass[11pt]{article}

% ---------- Packages ----------
\usepackage[margin=1in]{geometry}
\usepackage{amsmath, amssymb, amsthm}
\usepackage{mathtools}
\usepackage{enumitem}
\usepackage{fancyhdr}
\usepackage{xcolor}
\usepackage{hyperref}
\hypersetup{colorlinks=true, linkcolor=blue, citecolor=blue, urlcolor=blue}

% ---------- Header / Footer ----------
\setlength{\headheight}{14pt}
\pagestyle{fancy}
\fancyhf{}
\lhead{MA521 --- Symmetry, Geometry, and Representations}
\rhead{Homework 1}
\cfoot{\thepage}
\renewcommand{\headrulewidth}{0.4pt}

% ---------- Theorem-like environments ----------
\theoremstyle{plain}
\newtheorem{theorem}{Theorem}[section]
\newtheorem{lemma}[theorem]{Lemma}
\newtheorem{corollary}[theorem]{Corollary}
\newtheorem{proposition}[theorem]{Proposition}

\theoremstyle{definition}
\newtheorem{definition}[theorem]{Definition}

% ---------- Problem / Solution environment ----------
% \begin{problem}[optional label] ... \end{problem} typesets the given
% statement; \begin{solution} ... \end{solution} typesets the write-up
% and appends a QED-style marker at the end.
\newcounter{problemnum}
\newenvironment{problem}[1][]{%
  \refstepcounter{problemnum}%
  \par\medskip\noindent\textbf{Problem \theproblemnum{#1}.}\ \itshape
}{\par\medskip}

\newenvironment{solution}{%
  \par\noindent\textit{Solution.}\
}{%
  \hfill$\blacksquare$\par\bigskip
}

% ---------- Common shorthands ----------
\newcommand{\Hom}{\operatorname{Hom}}
\newcommand{\Aut}{\operatorname{Aut}}
\newcommand{\End}{\operatorname{End}}
\newcommand{\Iso}{\operatorname{Iso}}
\newcommand{\Span}{\operatorname{span}}
\newcommand{\Dim}{\operatorname{dim}}
\newcommand{\R}{\mathbb{R}}
\newcommand{\Q}{\mathbb{Q}}
\newcommand{\Z}{\mathbb{Z}}
\newcommand{\C}{\mathbb{C}}
\newcommand{\F}{\mathbb{F}}
\newcommand{\dual}[1]{#1^{\ast}}

\title{\vspace{-2em}MA521: Symmetry, Geometry, and Representations
\\ \large Homework 1}
\author{Asher Lucas-Cuddeback}
\date{04/09/2026}

\begin{document}
\maketitle
\thispagestyle{fancy}

\begin{problem}
  \
  \begin{enumerate}[label=(\alph*)]
    \item If $G$ and $H$ are two groups, explain why the set of group
      homomorphisms
      \[
        \Hom(G, H)
      \]
      only carries a natural group structure when $H$ is Abelian, and
      in that case, why $\Hom(G,H)$ becomes an Abelian group.

    \item If $V$ and $W$ are two $F$-vector spaces, show that the set
      of linear maps
      \[
        \Hom_F(V, W)
      \]
      naturally forms an $F$-vector space.

    \item Show that if $X$ is any mathematical object (set, group,
      ring, vector space, etc.), then the set of automorphisms
      $\Aut(X)$ naturally forms a group under composition.

    \item Show that if $A$ is any Abelian group, then the set of
      endomorphisms $\End(A)$ naturally forms a ring. (What natural
      operations of addition and multiplication would $\End(A)$ carry?)
  \end{enumerate}
\end{problem}

\begin{solution}
  \begin{proof}[Proof of 1(a)]
    Let \( f,g\in \Hom(G,H) \). For the group operation, select
    the group operation within the image space such that, \( \forall a\in G\):
    \[
      (f\cdot g)(a)=f(a)*g(a).
    \]
    With the composition operation selected, we must now show that
    the composed \( f\cdot g \) is itself a homomorphism. For this,
    we must have, \( \forall a,b\in G \):
    \[
      (f\cdot g)(ab)=(f\cdot g)(a) * (f\cdot g)(b)
    \]
    We expand both sides
    \[
      f(ab)*g(ab)=f(a)*g(a)*f(b)*g(b)
    \]
    \[
      f(a)*f(b)*g(a)*g(b)=f(a)*g(a)*f(b)*g(b)
    \]
    At this point, we are seemingly at an impasse. However, if we assert that
    \( H \) is an abelian group and recognize \( f(x) \) as an
    element in \( H \),
    we can 'slide' the problematic terms past each other to obtain:
    \[
      f(a)*g(a)*f(b)*g(b)=f(a)*g(a)*f(b)*g(b)
    \]

  \end{proof}

  \begin{proof}[Proof of 1(b)]
    For \( \Hom_F(V, W) \) to be a natural vector space, then \(
    \forall v\in V \) we must be able to define
    a vector addition and scalar multiplication operation which
    respects the vector space definitions. Let \( f, g \in
    \Hom_F(V,W) \) and \( \lambda \in F \), define
    \[
      (f + g)(v) = f(v) + g(v) \\
      \text{and} \\
      (\lambda f)(v) = \lambda f(v)
    \]
    The map \( (f + g) \) is linear since \( \forall v, w \in V \),
    \( \lambda\in F \)
    \begin{align*}
      (\lambda f + g)(v) &= (\lambda f)(v) + g(v) \\
      &= \lambda f(v) + g(v)
    \end{align*}
    thus \( (\lambda f + g) \in \Hom_F(V, W) \) and \( \Hom_F(V,W) \)
    is a vector space.
  \end{proof}

  \begin{proof}[Proof of 1(c)]
    The set of automorphisms is defined as \( \Aut (X) = \Iso(X, X)
    \). As a consequence, for any object \( x\in X \) and
    automorphism, \( f\in\Aut(X) \) we have:
    \[
      f(x)\in X
    \]
    Which allows us to select function composition as the composition
    operation in \( \Aut(X) \). This means that \( f,g\in \Aut(X) \)
    \[
      (f\cdot g)(x) = f(g(x)).
    \]
    We can now go through the checks to validate the group nature
    \begin{enumerate}
      \item Identity: Let \( 1\in \Aut(X) \) with \( 1(x) = x \)
      \item Closure: We get this property for free since all elements
        are automorphisms
      \item Inverses: since \( f\in\Aut(X) \), by definition, \(
        \exists f^{-1} \in \Aut(X) \) such that \( f^{-1}(f(x)) = x \)
      \item Associativity: let \( f,g,h \in \Aut(X) \).
        \begin{align*}
          ((f \cdot g) \cdot h)(x) &= (f\cdot g) (h(x)) \\
          &= f(g(h(x)))
        \end{align*}
        and
        \begin{align*}
          (f\cdot (g \cdot h))(x) &= f \cdot g(h(x)) \\
          &= f(g(h(x)))
        \end{align*}
    \end{enumerate}
    And so \( \Aut(X) \) is a group under function composition.
  \end{proof}

  \begin{proof}[Proof of 1(d)]
    Recall one of the 'cheating' ways to define a ring: a ring is a
    set, \( R \) equipped with 'addition' and 'multiplication'
    operations which interact via the distributive law. These
    operations must meet the following criteria
    \begin{enumerate}
      \item \((R, +)\) forms an abelian group.
      \item \( (R, \cdot) \) forms a semigroup (a group without inverses)
      \item \( \forall a,b,c,d\in R \), \( (a+b)\cdot(c+d)= ac + ad + bc + cd \)
    \end{enumerate}
    For criterion 1. we can reuse the result from part (a). Let \(
    (A, +) \) be an abelian group under \( + \). \( \End(A)=\Hom(A,A)
    \) so \( \forall a \in A \), \( f\in\End(A) \) we have \( f(a)\in
    A \). Thus, \( \End(A) \) inherits the abelian nature of \( A \) if
    \[
      (f + g)(a) = f(a) + g(a)
    \]
    since \( f(a),g(a)\in A \) and \( f,g\in\Hom(A,A) \). For
    criterion 2. we select the 'multiplication' operation to be
    function composition. We can again reuse results since we have
    already proved that function composition is associative in part
    (c). That only leaves closure and identity. Closure is easy since
    \( A \) is already closed:
    \[
      (f\cdot g)(a) = f(g(a)).
    \]
    And since \( g(a)\in A \), \( f(g(a))\in A \). As for identity,
    define the identity element to be the same as in part (c).
    \[
      (1\cdot f)(a) = f(a).
    \]
    Finally, we must show that the distributive law holds. Let \(
    i,j,k,l\in \End(A) \) and \( a \in A \).
    \begin{align*}
      ((i + j) \cdot (k + l))(a) &= (i + j)(k(a) + l(a)) \\
      &= i(k(a) + l(a)) + j(k(a) + l(a)) \\
      &= i(k(a)) + i(l(a)) + j(k(a)) + j(l(a)) \\
      &= (i\cdot k)(a) + (i\cdot l)(a) + (j\cdot k)(a) + (j\cdot l)(a) \\
      &= (i\cdot k + i\cdot l + j\cdot k + j\cdot l)(a)
    \end{align*}
    which is the desired result, and we are done.
  \end{proof}

\end{solution}

\begin{problem}
  Let $V$ be a finite dimensional $F$-vector space. Recall that the
  dual space of $V$ is the space: \[\dual{V} = \Hom_F(V, F).\]
  Explain why \((V^*)^*\) is naturally isomorphic to $V$.
\end{problem}

\begin{solution}
  Let \( v\in V \) and \( f\in V^* \). Given that \( V^*\cong \Hom(V,
  F) \), we know that \( f \) must respect the structure of \( V \).
  Define the map
  \[
    \begin{matrix}
      V \times \dual{V} \to V^{**} \\
      (v, f) \mapsto f(v)
    \end{matrix}
  \]
  Thus,
  \[
    \begin{matrix}
      v\mapsto(f\mapsto f(v)) \\
      V\to V^{ ** }
    \end{matrix}
  \]
  which we must now simply show that this map is invertible for it to
  be an isomorphism. Unfortunately, at this point we must acknowledge
  that we can choose a basis for \( V \). Let \( \dim_F(V) = n\)
  \[
    \begin{matrix}
      V=\Span_F(e_1, \dots, e_n) \\
      V^*=\Span_F(\epsilon_1, \dots, \epsilon_n)
    \end{matrix}
  \]
  Since we select
  \[
    \epsilon_i(e_j)=
    \begin{cases}1 & \text{if } i=j \\ 0 & \text{if } i\neq j
    \end{cases}
  \]
  Since the dual space is a vector space, the dual of the dual simply
  consists of the maps which re-associates the scalar 1
  with the basis vectors in \( V \). Thus, \(
  \Dim_F(V^{**})=n=\Dim_F(V) \) so \( V\cong V^{**} \).
\end{solution}

\end{document}
